% !TEX root = ../main.tex

% \chapter{Calculus of Inductive Constructions in Lean 4}
\chapter{Calcolo delle costruzioni induttive}
\label{chap:cic}

\section{Motivazioni}
Per poter parlare del kernel di Lean 4, è necessario approfondire il modello logico che i progetisti del linguaggio hanno deciso di implementare
e dei motivi per è stato necessario estendere il calcolo delle costruzioni. I sistemi di tipo e linguaggi finora esaminati appartenenti al lambda
cubo (\ref{tab:barendregt_combinations}), sono ottimali per indagare le proprita e i limiti teorici, ma possiedono delle criticita nel'utilizzo e 
nel'implementazione concreta. Per ovviare a tali criticita si definiscono delle estensioni che non alterano l'espressivita, cosi da poter mantenere
i risultati precedentemente dimostrati. Sebbene tali problemi siano condivisi da tutti i vertici del lambda cubo, per semplificare la trattazione
si utilizzera il calcolo delle costruzioni (\ref{eq:coc_grammar}). La metodologia utilizzata per ovviare a tali svantaggi potra essere generalizzata
e adattata a sistemi meno espressivi, mantenendo i vantaggi ottenuti.

\subsection{Limiti delle codifiche alla Church}
Nel calcolo delle costruzioni (\ref{sec:coc}) non esisteno tipi di dato primitivi. Per rappresentare strutture come i numeri naturali, si
utilizzavano le \emph{codifiche alla Church}. In tale paradgma, i dati non sono entita statiche, ma funzioni di ordine superiore che modellano
il proprio schema di iterazione. Tali definizioni comportano principalmente l'inefficienza algoritmica e la difficolta della definizione di un 
principio di induzione condiviso tra tutti i tipi.

\subsubsection{Inefficienza algoritmica dei numeri naturali}
Il limite più evidente delle codifiche alla Church è la degradazione delle performance computazionali. Poiché un numero $n$ è codificato come 
l'applicazione $n$-esima di una funzione, operazioni che intuitivamente dovrebbero essere istantanee diventano costose.

\begin{itemize}
    \item 
        \textbf{Il predecessore di Kelly-Church:} 
        Calcolare il predecessore richiede una ricostruzione laboriosa della sequenza. L'algoritmo standard ha una complessità 
        temporale di $O(n)$, rendendo operazioni come la sottrazione o la divisione proibitive per valori grandi.
    \item 
        \textbf{Costo della normalizzazione e conversione:}
        Il controllo dell'uguaglianza definizionale ($\equiv$) è l'operazione più frequente eseguita dal kernel di Lean. Nelle codifiche 
        alla Church, i numeri non sono dati statici ma programmi; pertanto, per verificare se due termini sono uguali, il kernel è costretto 
        a "eseguire" interamente le funzioni coinvolte, riducendo alberi di applicazioni $\lambda$ potenzialmente enormi. Questo processo 
        satura rapidamente le risorse di memoria e rallenta il \emph{type checking}.
\end{itemize}

\subsubsection{Svantaggi logici}
Il limite piu invalidante per un sistema cone Lean delle codifiche alla Church è di di natura fondazionale. Nel calcolo delle costruzioni, tali
codifiche non garantiscono intrinsecamente il \emph{principio di induzione} nel senso dell'aritmetica.

\begin{itemize}
    \item 
        \textbf{Mancanza di dipendenza:} 
        La codifica di Church è intrinsecamente polimorfica ma fatica a integrarsi con la logica dipendente. Risulta estremamente complesso 
        formulare predicati che dipendano dal valore specifico del dato (ad esempio, proprietà aritmetiche come la commutatività) senza ricorrere 
        ad assiomi esterni che appesantiscono la coerenza del sistema.
    \item 
        \textbf{Difficoltà nella distinzione dei costruttori:} 
        Nel paradigma funzionale, non esiste un meccanismo nativo per garantire che rappresentazioni diverse corrispondano a entità distinte. 
        Dimostrare proprietà fondamentali come l'iniettività del successore o la disuguaglianza tra lo zero e il successore ($\text{zero} \neq 
        \text{succ}(n)$) richiedea costruzione manuale di funzioni discriminanti ad differenti per ogni tipo di dato, rendendo la logica del 
        sistema meno intuitiva.

    \item 
        \textbf{Problemi di estensionalità e uguaglianza:} 
        Poiché i dati sono codificati come funzioni, il sistema eredita i limiti della conversione funzionale $\eta$-riduzione. Due rappresentazioni 
        dello stesso valore che differiscono anche minimamente nella loro struttura logica interna (pur essendo estensionalmente identiche nel comportamento) 
        non vengono riconosciute come definizionalmente uguali ($\equiv$) dal kernel. Questa frammentazione impedisce al sistema di convergere verso una 
        forma canonica univoca per ogni dato, costringendo l'utente a gestire manualmente complesse prove di equivalenza proposizionale anche per identità 
        che, in una rappresentazione strutturale, risulterebbero ovvie per semplice calcolo.
\end{itemize}